\paragraph{Motivation for learning this}
In the previous set of lectures, we briefly mentioned \textbf{Langevin dynamics}. Given a target distribution $\pi$, you can initialize your samples anywhere $X_0$ then time-evolve them according to
\begin{align}
	X_{t+h} & = X_t + h \nabla \log \pi(X_t)  + \sqrt{2h} \, Z & \text{Discrete Time}\\
	dX_t & = \nabla \log \pi(X_t) \, dt + \sqrt 2 \, dW_t & \text{Continuous Time}
\end{align}
and as $t \to \infty$, your samples will converge to $X_\infty \sim \pi$ (target measure). 
\\

Langevin fell under a special class of dynamics called \textbf{Markov Processes / Markov Chains}. These lectures will walk us through some of the tools \& techniques to analyze these processes. Markov chains show us in many disciplines, so to give the reader a consistent viewpoint, I'll adopt my POV (which is a computational statistician interested in the math; meaning you only have oracle access to $\pi$, and sometimes only have oracle access to the normalize density $\tilde \pi$ associated with $\pi = \tilde \pi / Z$.), but hopefully the lessons can give the reader enough to generalize this knowledge to other problems.  
\\
\\
From this view point, there's a couple questions that jump out to me
\begin{enumerate}
	\item Can we develop other algorithms which can generate samples from $\pi$?
	\item Once we specify a sampling algorithm, can we quantify how quickly samples will converge to $\pi$?
	\item Many algorithms are easy to specify in continuous time, but we implement only in discrete time. Can we quantify how the discretization affects the analysis?
\end{enumerate}
To keep us grounded, we'll work through these questions by looking back to Langevin, but the theory will stay general.


\section{Discrete-Time Markov Chains}
Consider a stochastic process $(X_t)_t$, and it's corresponding marginal over time $\text{Law}(X_t) := \pi_t$. The name of the game is to track how a state $\pi_t$ evolves over time. In these lecture notes, we'll look at a specific class of dynamics, called Markov Chains
\begin{definition}
	[Time-Homogenous Markov Chain] A discrete-time stochastic process $(X_t)_t$ is a Markov Chain if $p(X_{t+1} | X_t, ..., X_0) = p(X_{t+1} | X_t)$, denoted as the transition kernel. It is time-homogenous because the kernel is time-independent.
\end{definition}
In this case, the dynamics of the system are completely governed by the transition kernel. 
\begin{itemize}
	\item When the state space $\Omega$ is continuous, $\pi$ is a probability density function ($\pi \geq 0$ and $\int_\Omega \pi(x) \, dx = 1$). The dynamics are
	\begin{align}
		\pi_{t+1}(y) = \int P(x,y)\,  \pi_t(x) \, dx, \qquad P(x,y) = p(X_{t+1} = y| X_t = x)
	\end{align}
	The notation of $P(x,y)$ is chosen s.t. it can be read, the probability of starting at $x$ and ending up at $y$.
	\item If the state space is discrete, then $\pi_t$ is a vector in the probability simplex ($\pi_t$ is a vector, $[\pi_t]_i \geq 0$ and $\sum_i [\pi_t]_i = 1$). The dynamics turn into a matrix equation
	\begin{align}
		\pi_{t+1} = \pi_t P, \qquad P_{ij} = \mathbb P(X_{t+1} = j | X_t = i)
	\end{align}
	Because $P_{ij}$ represents probabilities $\sum_j P_{ij} = 1$. This is called a \textbf{stochastic matrix}
\end{itemize}
Some other fun facts,
\begin{itemize}
	\item When you apply $P$ multiple times, you're evolving multiple steps. E.g. $P^k$ time evolves for $k$ steps.
	\item The eigenvector $\pi = \pi P$ corresponds to the \textbf{stationary distribution} / fixed point.
\end{itemize}

\begin{eg}
	Looking back at our motivating example, Langevin dynamics. It is also a Markov Chain
	\begin{align}
		dX_t &=  \nabla \log \pi(X_t) dt + \sqrt 2 dW_t\\
		X_{t + h} & = X_t + h\, \nabla \log \pi(X_t) + \sqrt{2h} Z & \text{(Discretization)}
	\end{align}
	First, let's prove it's Markovian
	
	Now let's derive the transition kernel $p(X_{t+h} | X_t = x)$
	\begin{align}
		X_{t + h} & = X_t + h\, \nabla \log \pi(X_t) + \sqrt{2h} Z\\
		X_{t + h} | X_t =x & \overset{d}{=} x + h \nabla \log \pi (x) +\sqrt{2 h} Z
	\end{align}
	The probability distribution of $Z$ is $\mathcal N(0, \mathbb I)$. We can use the rules of Gaussian random variables
	\begin{align}
		b + \sigma Z & \sim \mathcal N(b, \sigma^2 \mathbb I)\\
		\therefore x + h \nabla \log \pi(x) + \sqrt{2 h} & \sim \mathcal N( x + h \nabla \log \pi(x), \; 2h \mathbb I)
	\end{align} 
	So if the RHS is drawn from a Gaussian, therefore the LHS is also a Gaussian. Therefore the probability distribution of $X_{t+h} = x' | X_t=x$ is
	\begin{align}
		p(x' | x) = \frac{1}{Z} \exp \Big( (x' - (x + h \nabla \log \pi(x)))^2 / 4 h \Big)
	\end{align}
\end{eg}


\subsection{Stationary Distribution}
Since our objective is to construct a Markov chains which has a stationary distribution. It turns out a chain has a stationary distribution under two sufficient conditions
\begin{enumerate}
	\item \emph{Irreducibility:} It is possible to get from $x$ to another state $x'$ w/ probability $>0$ in a finite number of steps
	\item \emph{Aperiodicity:} It is possible to return to any state at any time. I.e. $\exists n$ s.t. $\forall \omega \in \Omega$ and $\forall n ' > n$, $\mathbb P(X_{n'} = i | X_0 = i) >0$.
\end{enumerate}
The first prevents absorbing states, e.g. states that we can never leave. The second prevents transition operations such as
\begin{align}
	T = \begin{pmatrix}
		0 & 1 \\ 1 & 0
	\end{pmatrix}
\end{align}
which is a chain that alternates between 1 and 2 without ever settling. These conditions constrain the types of Markov chains to get a nice result.
\begin{theorem}
	[Perron-Frobenius] An irreducible \& aperiodic Markov Chain $P$ 
	\begin{enumerate}
		\item Admits a unique stationary distribution $\pi$
		\item The stationary distribution $\pi$ is the eigenvector of $P$ corresponding to the largest eigenvalue $|\lambda| = 1$.
	\end{enumerate}
\end{theorem}
The proof is more about linear algebra than Markov Chains, so it's omitted.\\
\\
The importance of the theorem is that the eigenvector corresponding to the largest eigenvalue is the stationary distribution. Visualize the coordinate transformation of the matrix $P$. Repeated applications of $P$ will slowly project your vector along eigenvector
\begin{fact} For an irreducible \& aperiodic Markov Chain $P$ with stationary distribution $\pi$, repeated applications of $P$ will make any initialization $\pi_0$ converge to $\pi$.
	\begin{align}
		\lim_{k \to \infty} \pi_0 P^k = \pi
	\end{align}
\end{fact}
The proof will be shown in the section on convergence rates.


\paragraph{Detailed Balance and Metropolis-Hastings Algorithm}
You might wonder, given oracle access to $\pi$, how can we construct a Markov chain $P$ which emits $\pi$ as the stationary distribution. \\
\\
Well, say you're at the stationary distribution $\pi$, then you'd hope for the following properties
\begin{enumerate}
	\item Invariant under time-evolution: $\pi = \pi P$
	\item Invariant under reverse time-evolution $\pi_i = \sum_j p(x_{t+1} = j | x_{t} = i) \pi_j \implies \pi = P^T \pi$
\end{enumerate}

\begin{definition}
	[Detailed Balance] It is the balance between transition rates
	\begin{align}
		\pi(y) p(x | y) = \pi(x) p(y | x)
	\end{align}
	where $p$ is the transition kernel of a Markov Chain.
\end{definition}
\begin{fact}
	Let $p$ be the transition kernel associated with a Markov chain. If  $(p, \pi)$ satisfies detailed balance, then $\pi$ is the stationary distribution of the chain.
\end{fact}
\begin{proof}
	\begin{align}
		\pi(y) p(x|y) & = \pi(x) p(y|x)
	\end{align}
	Just integrate w.r.t. $y$	\begin{align}
		\sum_y \pi(y) p(x|y) & = \pi(x) \sum_y p(y|x)\\
		\sum_y \pi(y) p(x|y) & = \pi(x)& \sum_y p(y| x) = 1\\
		\pi P & = \pi 
	\end{align}
	and then $x$
	\begin{align}
		\pi(y) \sum_x p(x| y) & = \sum_x \pi(x) p(y|x)\\
		\pi& = \pi P
	\end{align}
	"It seems too easy", Jonathan Goodman.
\end{proof}
As the person running experiments, I simply want to specify a target distribution $\pi$ \& a proposal kernel $p$, and check which one works well. However not every transition kernel satisfies detailed balance for arbitrary $\pi$. Perhaps we can create a framework which always emits detailed balance.\\
\\
Consider 




Metroplist-Hasting rejection is a way to create a transition kernel $p(\cdot | x)$ which satisies detailed balance.
\begin{align}
	\alpha(x | x') = \min \left( \frac{\pi(x) p(x' | x)}{\pi(x') p(x | x')}, \ 1 \right)
\end{align}

\begin{algorithm}
\caption{Metropolis--Hastings}
\begin{algorithmic}[1]
\Require Target $\pi$, proposal $p(\cdot\mid x)$, initial $X_0$, steps $N$
\For{$t = 0, 1, \dots, N-1$}
    \State Propose $X' \sim p(\cdot \mid X_t)$
    \State $\alpha \gets \min\!\left(1,\ \dfrac{\pi(X')\,p(X_t \mid X')}{\pi(X_t)\,p(X' \mid X_t)}\right)$
    \State Draw $u \sim \mathrm{Uniform}(0,1)$
    \If{$u \le \alpha$}
        \State $X_{t+1} \gets X'$ \Comment{accept}
    \Else
        \State $X_{t+1} \gets X_t$ \Comment{reject}
    \EndIf
\EndFor
\State \Return $(X_0, \dots, X_N)$
\end{algorithmic}
\end{algorithm}

\begin{eg}
	[Metropolis-Adjusted Langevin (MALA)] For Langevin dynamics, we proved that it has the stationary distribution $\pi$. However this was in continuous time. Does this still hold when we discretized time?
	\\
	\\\
	Recall the Euler-Maruyama discretized Langevin
	\begin{align}
		X_{t+h} & = X_t + h \nabla \log \pi(X_t) + \sqrt{2h} \, Z\\
		p(x' | x) & = \frac{1}{Z} \exp \left( -\frac{(x' - x - h \nabla \log \pi(x))^2}{4h}\right)
	\end{align}
	this is known as \textbf{Unadjusted Langevin (ULA)}. It's clear the kernel doesn't satisfy detailed balance (for arbitrary $\pi$), as there's no $\pi$ term in the kernel. Although detailed balance is "sufficient, not necessary" that you converge to $\pi$, turns out the discretization makes you sample a biased version of $\pi$. We can correct for this bias by using the proposal kernel in the Metropolis-Hasting step.\\
	\\
	Let's calculate the acceptance ratio for the ULA kernel. This procedure is called \textbf{Metropolis-Adjusted Langevin (MALA)}
	\begin{align}
		\log \frac{p(x'| x)}{p(x|x')}  & = \frac{1}{4h} \left[ (x' - x + h \nabla \log \pi(x))^2 - (x - x' + h \nabla \log \pi(x'))^2 \right]\\
		& = \frac{h}{4} (\nabla \log \pi(x)^2 - \nabla \log \pi(x')^2)  + \frac{1}{2}\langle x' - x, \nabla \log \pi(x) + \nabla \log \pi(x')\rangle
	\end{align} 
	which you plug into $\alpha = \min \left( \frac{\pi(x) p(x' | x)}{\pi(x') p(x | x')}, \ 1 \right)$.\\
	\\
	You can also run the Langevin 
\end{eg}
\begin{problem}
	The previous MALA algorithm was stated for single step. Consider a modification to the MALA algorithm for multiple steps
	
	What is the new correction factor $p(x' | x)/p(x|x')$ in the Metropolis acceptance ratio?
\end{problem}

\paragraph{Sampling without Detailed Balance} 
For example: sequential monte carlo \& diffusion models, bouncy particle sampler.


\subsection{Convergence to Stationary Distribution}
It's nice that we can construct Markov chains where we know the stationary distribution. However, I need to run this and get some results, so perhaps we can engineer transition kernels with quicker convergence?

\paragraph{Discrete Space \& Discrete Time Markov Chains}
Let's now prove convergence fact from the previous section, the proof will give us a nice insight into convergence rates.
\begin{fact}
	Consider an irreducible \& aperiodic Markov Chain $P$ with stationary distribution $\pi$. Repeated applications of $P$ will make any initialization $\pi_0$ converge to $\pi$.
	\begin{align}
		\lim_{k \to \infty} \pi_0 P^k = \pi
	\end{align}
\end{fact}
\begin{proof}
	For simplicity (and w/o loss of generality), assume $P$ is diagonalizable with eigensystem $\{(\lambda_i, v_i)\}_{i=1}^d$, where $\lambda_1 > ... > \lambda_d$. By Perron-Frobenius theorem $\lambda_1 = 1$ and $v_1 = \pi$.\\
	\\
	Decompose $\pi_0$ into a linear combination of the eigensystem
	\begin{align}
		\pi_0 P^k & = (c_1 \pi + c_{2} v_{2} ... + c_d v_d) P^k\\
		& = c_1 (1)^k \pi+ c_{2} \lambda_{2}^k v_{2} + ... + c_d \lambda_d^k v_d\\
		\lim_{k \to \infty}\pi_0 P^k & =  c_1 \pi + \mathcal O(\lambda_{2}^k) = c_1 \pi
	\end{align}
	You can kill off $\lambda_{2}$ (and below) because it's smaller than $\lambda_2 = 1$.
	Because the state still lives in $\mathcal P(\Omega)$, it must normalize, therefore $c_1 = 1$.\\
	\\
	Note: power iteration can be generically applied to find the eigenvector of the largest eigenvalue. The convergence rate is generically controlled by the \textbf{spectral gap} $|\lambda_1| - |\lambda_2|$.
\end{proof}
Notice, for finite $k$ you're converging to the stationary distribution in corrections of $\lambda_{2}$ (second largest eigenvalue). This means spectral properties of the matrix controls the convergence rate of our Markov Chains.
\begin{definition}
	[Spectral Gap] The spectral gap for a matrix $P$ is defined as
	\begin{align}
		\gamma = |\lambda_1| - |\lambda_2|
	\end{align}
	For irreducible \& aperiodic Markov chains, $\lambda_1 = 1$, meaning the convergence is entirely controlled by $\lambda_2$.
\end{definition}
\paragraph{Continuous Space \&  Discrete/Continuous Time Markov Chains}
When the state space is continuous, the time-evolution $P$ gets promoted to an operator. And you can imagine finding eigenvalues of some operator will be quite difficult. We need to change our tools a little bit!





\section{Continuous-Time Markov Chains}
When we go to continuous space Markov Chains, it's easy to start with continuous time and use the theory of ODE/PDEs to make statements, then bound the discretization error.
\begin{definition}
	[Markov Semigroup Operator] Given a Markov Process $(X_t)_{t\geq 0}$, tis semigroup $(P_t)_{t\geq0}$ acts on test functions $f : \mathbb R^d \to \mathbb R$ by
	\begin{align}
		(P_t f) (x) \equiv \mathbb E[f (X_t) | X_0 = x]
	\end{align}
\end{definition}
Why is it called a semigroup? In math a semigroup is when you have closure under composition $g_1 \circ g_2 \in G$, and have an identity $e \circ g = g \circ e \in G$ (s.t. $e \in G$). If you have closure under inverse, then this is promoted to be a group. 
\begin{property}
	Let's figure out in what sense this is a semigroup
	\begin{enumerate}
	\item Identity: $P_0 = \textsf{Id}$.
	
	\emph{Proof:} $(P_0 f) (x) = \mathbb E[f(X_0) | X_0 = x] = f(x)$, so $P_0$ acts as an identity operator.
	
	\item Composition: $P_t \circ P_s = P_{t+s}$.
	
	\emph{Proof:} 
	\begin{align}
		(P_s f)(x) & = \mathbb E[f(X_s) | X_0 = x]\\
		(P_t P_s f)(x) & = \mathbb E[\mathbb E[f(X_s) | X_0 = X_t] | X_0 = x]\\
		& = \mathbb E[\mathbb E[f(X_{s+t}) |X_t] | X_0 =x]\\
		& = \mathbb E[f(X_{s+t}) | X_0 = x] = (P_{t+s} f)(x)
	\end{align}
\end{enumerate}
\end{property}




\begin{definition}
	[Generator] The generator $\mathcal L$ is defined as
	\begin{align}
		\mathcal L f := \lim_{h \downarrow 0} \frac{P_h f - f}{h}
	\end{align}
\end{definition}

\begin{eg}
	Let's compute the infinitesimal generator for Langevin dynamics. Since we'll take the time-increment to be small, we can start at the Euler-Maruyama discretization.
	\begin{align}
		X_h = X_0 + h \nabla \log \pi(X_0) + \sqrt{2 h} Z
	\end{align} 
	where $Z \sim \mathcal N(0, \mathbb I)$. To compute $P_h f = \mathbb E[f(X_h) | X_0 = h]$, we can apply a test function $f$ to $X_h$, and then Taylor expand to second order, that way remaining terms are order $h$ (taking the limit $h \to 0$ in the infinitesimal generator will kill $h^{3/2}$ terms).
	\begin{align}
	f(X_h) & = f(X_0) + f'(X_0) (X_h-X_0) + \frac 1 2 f''(X_0)(X_h-X_0)^2 + \mathcal O((X_h-X_0)^3)\\
	& = f(X_0) + \langle \nabla f(X_0) ,  \, -h \nabla \log \pi(X_0)  + \sqrt{2h} Z\rangle \nonumber \\
	& \quad + \frac{1}{2} (-h \nabla \log \pi(X_0)  + \sqrt{2h} Z)^T \nabla^2 f(X_0) (-h \nabla \log \pi(X_0)  + \sqrt{2h} Z)  
\end{align}
Now we can take the expectation value
\begin{align}
	\mathbb E[f(X_h) & | X_0=x] = -h \langle \nabla f(x) , \nabla \log \pi(x) \rangle +  h \, \text{Tr}[\nabla^2 f(x)]
\end{align}
Plugging into the definition of the generator
\begin{align}
	\mathcal L f =  \langle f , \nabla \log \pi \rangle + \Delta f
\end{align}
\end{eg}
\begin{problem}
	Find the infintesimal generator associated for a generic stochastic process
	\begin{align}
		dX_t= \mu_t(X_t) dt + \sigma_t dW_t
	\end{align}
	Does it look familiar?
\end{problem}

\subsection{Convergence to Stationary Distribution}
To quantify the convergence rate to the stationary distribution, we could inspect a "distance" measure between the marginal $p_t$ of the Markov Chain and the target $\pi$, and see if its growing / shrinking in time.
\begin{definition}
	[$f$-Divergence] Given a convex function $f: \mathbb R \to \mathbb R$, the $f$-Divergence is
	\begin{align}
		D_{f}(p \| \pi) \equiv \mathbb E_{p} \left[f  \left( \frac{p }{\pi} \right) \right] = \int p(x) \, f\left(\frac{p(x)}{\pi(x)}\right) \, dx
	\end{align}
	A common choice for $f := \log$; this choice is called the Kullback-Leibler (KL) Divergence
	\begin{align}
		\textsf{KL}(p \| \pi) := D_{\log} (p \| \pi) = \int p(x) \, \log \left(\frac{p(x)}{\pi(x)}\right) \, dx
	\end{align}
	It is preferred by computational staticians because it can be numerically estimated $\textsf{KL}(p \| \pi) = \mathbb E_p [\log p]   - \mathbb E_p [\log \pi]$
\end{definition}
\begin{property}
	\begin{enumerate}
		\item Not symmetric: $D_f(p \| \pi) \neq D_f(\pi \| p)$
		\item Positive-definite: $D_f(p \| \pi) \geq 0$, and $D_f(p \| q) = 0 \iff p = q$		
	\end{enumerate}
\end{property}
I am introducing this new family of divergences, instead of using $\textsf{TV}(p, \pi) = \frac{1}{2} \int |p(x) - \pi(x)| \, dx$ because the $\textsf{TV}$ is not differentiable. 
\begin{eg} Let's compute the $f$-Divergence's time evolution under Langevin. Recall the law for Langevin for a target measure $\pi \propto e^{-U(x)}$ is
	\begin{align}
		\partial_t p_t = - \nabla \cdot (\nabla U\, p_t) + \Delta p_t
	\end{align}
	\begin{align}
		\partial_t D_f(p_t \| \pi) & = \int \dot p_t (x) \, f \left(\frac{p_t(x)}{\pi(x)} \right) + \frac{p_t(x) \dot p_t(x)}{\pi(x)} f' \left(\frac{p_t(x)}{\pi(x)} \right)\\
		& = \int \dot p_t(x) \left[  f \left(r_t \right) + r_t f '\left(r_t\right) \right] & r_t \equiv p_t / \pi\\
	\end{align}
	If we specialize to the \textsf{KL} divergence, $f  := \log$
\end{eg}




\begin{definition}
	[Log-Sobolev Inequality (LSI)] $\pi$ satisfies a Log-Sobolev Inequality with constant $\lambda >0$, if
	\begin{align}
		\forall p \in \mathcal P, \quad \textsf{KL}(p \| \pi) \leq \frac{1}{2\lambda} \textsf{FI}(p \| \pi)
	\end{align}
\end{definition}


\begin{fact}
	[Bakery-Emery Criterion] Suppose that $\pi = e^{-U}$ with $U$ $\beta$-strongly convex and $\lambda$-Lipschitz, that is $\lambda \mathbb I \preceq \nabla^2 f \preceq \beta \mathbb I$. Then $\pi$ satisfies LSI with constant $\lambda$.
\end{fact}




\begin{thebibliography}{XYZ99}  % widest label, for alignment

\bibitem[Che26]{SinhoLC}
Sinho Chewi. \emph{Log Concave Sampling}. \url{https://chewisinho.github.io/main.pdf}.

\bibitem[KE25]{KuleshovErmonPGM}
Volodymyr Kuleshov and Stefano Ermon. \emph{Probabilistic Graphical Models Lecture Notes}. \url{https://ermongroup.github.io/cs228-notes/}

\bibitem[Bru25]{BrunaIR}
Joan Bruna. \emph{Inference and Representation Lecture Notes}. \url{https://joanbruna.notion.site/ir25}

\bibitem[Bru24]{BrunaMDL}
Joan Bruna. \emph{Mathematics of Deep Learning Lecture Notes}. \url{https://joanbruna.notion.site/Mathematics-of-DL-2c72610b1a9c401dbf624b3355fa0f2d}


\bibitem[Goo20]{Goodman}
Jonathan Goodman. \emph{Markov Chain Monte Carlo Lecture Notes}. \url{https://math.nyu.edu/~goodman/teaching/MonteCarlo20/MonteCarloMethods.html}

\bibitem[Owe13]{Owen}
Art B. Owen. \emph{Monte Carlo theory, methods, and examples}. \url{https://artowen.su.domains/mc/}
\end{thebibliography}





